% Vietnamese translation for OI Public Library
% Source-File: IOI中国国家候选队论文/2003/项荣璟/项荣璟.ppt
\documentclass[11pt]{article}
\usepackage{oipl-vn}

\title{Bản trình chiếu: Quy hoạch động chia đoạn}
\author{Xiang Rongjing}
\date{2003}

\begin{document}
\maketitle

\origtitle{Slide companion for interval dynamic programming}
\sourcepath{IOI China candidate team papers / 2003 / Xiang Rongjing / Xiang Rongjing.ppt}

\section{Phạm vi bản dịch}

Tệp gốc chỉ có bản PowerPoint. Phần chữ khôi phục được chứa tên \textit{CERC
1998}, \textit{IOI 2002}, một ví dụ chia dãy \(100,200,\ldots,900\), và các
công thức \(f(i,j)\), \(g(i,k)\), \(s_{i,k}\). Điều này cho thấy trọng tâm là
quy hoạch động trên đoạn hoặc phân hoạch dãy.

\section{Mô hình}

Ta xét một dãy cần chia thành các đoạn liên tiếp. Chi phí của một đoạn có thể
phụ thuộc vào tổng \(T_i+\cdots+T_j\), thời gian chuẩn bị \(s\), hoặc các giá
trị \(F_i\). Ví dụ trong slide có phân hoạch
\[
  \{1,2\},\{3\},\{4,5\}.
\]

\section{Chuyển trạng thái}

Các ký hiệu khôi phục được cho thấy hai bảng chính:
\[
  f(i,j),\qquad g(i,k).
\]
Một dạng chuyển thường gặp là tách đoạn \([i,n]\) tại vị trí \(j\), rồi kết
hợp lời giải của \([i,j]\) với \([j+1,n]\). Các chỉ số \(s_{i,k}\) gợi ý tối
ưu hóa điểm chia khi tính \(g(i,k)\).

\section{Thông điệp}

Bản trình chiếu minh họa cách biến một bài phân hoạch thành công thức DP và
sau đó giảm độ phức tạp từ \(O(n^3)\) xuống \(O(n^2)\) bằng cách khai thác
tính đơn điệu của điểm chia.

\end{document}
